\documentclass[a4paper, 11pt]{article}
\usepackage[top=2.5cm, bottom=2.5cm, left=2.5cm, right=2.5cm]{geometry}
\usepackage[utf8]{inputenc}
\usepackage[T1]{fontenc}
\usepackage[italian]{babel}
\usepackage{amsmath, amssymb}
\usepackage{siunitx}
\usepackage{setspace}
\usepackage{titlesec}
\usepackage[colorlinks=true, urlcolor=blue, linkcolor=black]{hyperref}

\setstretch{1.15}
\setlength{\parindent}{0pt}
\setlength{\parskip}{0.5em}


\titleformat{\section}
  {\normalfont\large\bfseries}
  {\thesection}{1em}{}
\titlespacing*{\section}{0pt}{1.5ex plus .5ex minus .2ex}{0.8ex plus .1ex}

\pagestyle{plain}

\begin{document}


\section*{Sintesi dell'elaborato}


La microscopia a tempi di vita di fluorescenza (FLIM, FLuorescence Lifetime Imaging Microscopy), è una tecnica spettroscopica che fornisce informazioni sulle specie fluorofore presenti in un campione in modo indipendente dalla loro concentrazione.

L'analisi dei dati FLIM si basa tradizionalmente su procedure di fitting multi-esponenziale del decadimento esponenziale per ciascun pixel dell'immagine, e risulta computazionalmente costosa e sensibile al rumore. Un approccio più moderno è invece quello dell'analisi fasoriare: il decadimento temporale di ciascun pixel viene proiettato in un punto dello spazio complesso $\mathcal{P}(\omega) = g(\omega) + i\,s(\omega)$ attraverso una trasformata di Fourier calcolata alla frequenza angolare di modulazione del sistema ($\omega =2\pi/T_{laser}$), e normalizzata dal numero totale di fotoni raccolti nel pixel.\\
\[
G(\omega) = \frac{\int_0^{T} I(t) \cos(\omega t)\,dt}{\int_0^{T} I(t)\,dt}, \qquad 
S(\omega) = \frac{\int_0^{T} I(t) \sin(\omega t)\,dt}{\int_0^{T} I(t)\,dt}
\]

In questo spazio, punti associati ai decadimenti mono-esponenziali ($I(t) = I_0 e^{-t/\tau}$) cadono sul \textit{semicerchio universale} (centrato in $(G, S)=(0.5, 0)$ e a valori nel primo quadrante), mentre i multi-esponenziali (emissione di specie fluorofore con tempi di vita diversi) si trovano nell'area compresa tra questa curva di riferimento e l'asse reale. In quest'ultimo caso, grazie alla linearità della trasformata di Fourier, i fasori obbediscono alla regola della somma vettoriale pesata sull'intensità frazionaria dei singoli fluorofori.\\

Nei dati FLIM reali, i decadimenti di fluorescenza sono colvoluti temporalmente con una funzione di risposta dello strumento (IRF, \textit{Instrument Response Function}), fortemente dipendente dalla strumentazione utilizzata. Nell'analisi fasoriale, è possibile sfruttare il teorema della convoluzione per isolare il segnale di fluorescenza dalla IRF mediante una semplice divisione complessa sul piano dei fasori: $z_{\text{decadimento}} = z_{\text{segnale}}/z_{\text{IRF}}$.

Nella prima fase operativa è stato sviluppato un modello in linguaggio Python per la generazione e l'analisi di segnali di fluorescenza simulati, volto al collaudo delle funzioni di analisi prima dell'applicazioni a raccolte di dati reali.\\
In primo luogo sono stati simulati decadimenti prima mono-esponenziali, poi multi esponenziali (a due e tre componenti), al variare di $\tau_k$, $f_k$ (rispettivamente i tempi di vita e le intensità frazionarie di ogni componente) e del rapporto segnale-rumore. Dopodiché i segnali sono stati rappresentati sul piano dei fasori. Per i decadimenti a doppia esponenziale è stata confermata visivamente la posizione del fasore sul segmento congiungente le due componenti pure; le rispettive frazioni di intensità risultano inversamente proporzionali alle distanze dai vertici ($f_A : f_B = \overline{PB} : \overline{PA}$).\\
Analogamente, per i decadimenti tri-esponenziali il fasore $P$ cade all'interno del triangolo di vertici $A$, $B$ e $C$ (associati alle tre specie pure), e la frazione di intensità di ciascuna componente risulta proporzionale all'area del sotto-triangolo opposto al rispettivo vertice.

Successivamente è stata modellizzata una IRF gaussiana a $\mu$ e $\sigma$ variabili, ed è stata aggiunta a segnali mono-esponenziali e bi-esponenziali mediante una convoluzione sul piano temporale. Segnale totale e deconvoluto, IRF e componenti pure deconvolute sono stati rappresentati sul piano dei fasori, confermando il funzionamento dell'algoritmo di deconvoluzione implementato sul piano dei fasori.\\

Nella sezione principale della trattazione, gli argoritmi di analisi sviluppati in precedenza sono stati applicati a immagini FLIM reali, ovvero dataset contenuti in file \textit{tiff} di dimensioni \texttt{x}, \texttt{y} e \texttt{time}, ovvero in cui per ogni singolo pixel è contenuto il segnale del decadimento di fluorescenza.

Il primo campione biologico analizzato è un embrione di Zebrafish a 3 giorni dalla fecondazione, di cui viene studiata la fluorescenza endogena.\\
Il fasore della IRF è stato estratto grazie a un'immagine di calibrazione con fluorescina, $\tau = 4\text{ns}$, corredata all'immagine del campione.
Dopo la deconvoluzione dalla IRF, sono stati rappresentati i fasori associati ai singoli pixel, non prima di aver applicato un filtro median alla matrice delle coordinate $(G, S)$, per ridurre la dispersione dei fasori causata dal rumore statistico (o \textit{shot noise}). È inoltre stata applicata un soglia di esclusione dei pixel con conteggi di fotoni troppo bassi, e una procedura di selezione di regioni di interesse (ROI, Region Of Interest). Quest'ultima agisce come una maschera quadrata di lato e posizione variabile che può essere spostata sull'immagine per mappare sul piano dei fasori i solo i pixel selezionati.

Il secondo campione è un altro campione di Zebrafish ma a 5 giorni dalla fecondazione; in questo caso, oltre all'acquisizione di fluorescenza totale, sono state fornite anche due acquisizioni ottenute filtrando rispettivamente l'emissione nel rosso ($\lambda = 607\, \text{nm}$) e nel verde ($\lambda = 550\, \text{nm}$).\\
La calibrazione della IRF è stata effettuata sfruttando un segnale di generazione di seconda armonica (\textit{Second Harmonic Generation}, SHG), un processo istantaneo che vede la conversione di due fotoni infrarossi incidenti in un fotone visibile; successivamente tutti e 3 i dataset (totale, rosso e verde) sono stati rappresentati sul piano dei fasori.

Selezionate due ROI distinte sull'immagine, è stato determinato e messo a confronto il rapporto tra i contributi dei fluorofori rosso e verde in ciascuna regione, valutando la posizione del fasore medio di ogni ROI rispetto ai centroidi dei cluster rosso e verde. 
Infine, è stato effettuato uno studio delle incertezze valutando la matrice di covarianza di ciascun cluster e propagando l'errore alle frazioni di intensità tramite simulazione Monte Carlo. 
I risultati ottenuti mostrano, nella prima regione, una predominanza della componente rossa con frazioni pari a $f_{\text{rosso}} = 0.94 \pm 0.04$ e $f_{\text{verde}} = 0.06 \pm 0.04$, mentre nella seconda si osserva una presenza bilanciata di entrambe le componenti ($f_{\text{rosso}} = 0.55 \pm 0.12$ e $f_{\text{verde}} = 0.45 \pm 0.12$).\\

In conclusione, da questa trattazione emerge come l'analisi fasoriale possa fornire un'alternativa valida al \textit{fitting} non lineare, tagliando i costi computazionali di alcuni passaggi (interpolazione) e rendendo possibili alcune operazioni (deconvoluzione della IRF senza conoscerne la forma).\\
È stata infine sviluppata un'interfaccia grafica interattiva, implementata in una \textit{WebApp}, con lo scopo di agevolare la visualizzazione dei vari step sia della generazione di segnali, che della corrispondenza tra fasori e pixel di un'immagine.

\end{document}



